\documentclass[11pt,a4paper]{article}
\usepackage[margin=23mm]{geometry}
\usepackage{amsmath,amssymb}
\usepackage[hidelinks]{hyperref}
\setlength{\parindent}{0pt}
\setlength{\parskip}{6pt}
\raggedright
\begin{document}
{\Large BP-R Stage 3: interface continuation}\\[4pt]
7 October 2026 --- static review; S3-G2-Q discussion; Stage 3 open

\textbf{Status.} The previous partial result is accepted. Reference A,
upstream admission M.A and priority N.A remain closed. Neither compact route
Q.A nor Q.B is adopted. These are analytical checks, not traffic simulation,
parameter estimation, clearance validation or a gradient test.

\section*{1. Conservative count and event ownership}
Let $F_{cb}^k\ge0$ be actual continuous mass transferred from control volume
$c$ to $b$ during time bin $k$, $n_c^k$ its mass, and $I_c^k,E_c^k$ its external
arrival/exit mass. For road/ramp and entrance waiting volumes together,
\begin{equation}
n_c^{k+1}=n_c^k+I_c^k-E_c^k+
 \sum_b F_{bc}^k-\sum_b F_{cb}^k,
\qquad M^{k+1}-M^k=I^k-E^k.
\end{equation}
Internal transfers cancel. A descriptor adds no mass. The source's lateral
transfer advances to the next longitudinal target; it is not advected twice.
Receiving, sending and source step-size bounds remain separate requirements.

For stream $s$, $H_s(t)$ is actual cumulative transferred mass after the start,
$m_*=1$ vehicle-equivalent is the descriptor threshold, and
$z_s\in[0,1)$ is public initial threshold progress. At the stated source
sampling instants, the number of new descriptors is
\begin{equation}
K_s(t)=\left\lfloor z_s+H_s(t)/m_*\right\rfloor,
\qquad \sum_{s=1}^{S}K_s(t)\le LP/m_*+S.
\end{equation}
The bound is conditional on a finite acyclic directed mesh: every unit of the
finite offered/initial population $P$ crosses at most $L$ descriptor-generating
interfaces. $S$ is the finite number of such streams. Initial descriptors are
additional finite public state. The bound does not establish clearing or
justify continuous event-time interpolation of sampled source births.

For a source-free lane segment with cumulative-count field $N(x,t)$, density
$k=-N_x$, flow $q=N_t$, and constraint path $x_b(t)$, zero relative passing gives
\begin{equation}
\frac{d}{dt}N(x_b(t),t)=q(x_b,t)-\dot x_b\,k(x_b,t)=0.
\end{equation}
This is the local moving-boundary interface. Distributed lateral sources and
junction events require count anchoring/reindexing; equation (3) is not a
proof of their joined solution. Cell inventories, node counts and moving-path
values must be consistent. No source coincidence perturbation is adopted.

\newpage
\section*{2. A concrete dynamic compact source law}
Let $D_r,D_m$ be ramp and compatible mainline-group sending rates, $R_0$ the
\emph{ordinary, unmodified} target receiving rate, $Q$ target FD capacity,
$p\in(0,1)$ the approved fixed ramp priority, and $\delta\in[0,1)$ a fixed
phenomenological capacity-drop fraction. All flows have units veh-equivalent/s.
RM0010, Jin, Gan and Lebacque (2013), Section 6 equation (15), supplies
\begin{equation}
R_* = \min\{R_0,\ Q(1-\delta\,
 \mathbf 1_{\{D_r+D_m>R_0\}})\}.
\end{equation}
Its priority merge transcription is
\begin{equation}
f_r=\min\{D_r,\max\{R_*-D_m,pR_*\}\},\quad
f_m=\min\{D_m,\max\{R_*-D_r,(1-p)R_*\}\}.
\end{equation}
Use the already approved non-ramp subdivision. Densities evolve by equation
(1), with flows formed from the model's own state; approach admission remains
M.A. These wrapper/grouping adaptations are not the source's multilane theorem.
The source law is phenomenological, not a derived disturbance kinetic law.

Local bounds follow from the previously qualified priority operator:
\begin{equation}
0\le f_r\le D_r,\quad 0\le f_m\le D_m,\quad
f_r+f_m=\min\{D_r+D_m,R_*\}\le R_0.
\end{equation}
Thus this is an independent dynamic conservative candidate, rather than only a
mean-throughput formula. The trigger must use $R_0$, not already reduced $R_*$.
Bounds do not prove predictive adequacy, joined reference validity or gradients.

\textbf{Reviewer hand counterexample to the inactive match.}
Take no ramp entry or active moving constraint. Share the same current ordinary
density/receiving state across models and assign
\begin{equation}
D_r=0,\quad D_m=Q,\quad R_0=0.9Q,\quad\delta=0.2:
\qquad f_m^{\rm compact}=0.8Q<0.9Q=f_m^{\rm unobstructed}.
\end{equation}
These values are an algebraic witness, not calibrated data or a simulated
trajectory. The source model's own physical domain is not refuted. Its
$\delta=0$ limit gives the ordinary node; that different parameter limit does
not prove equality when constraints stop binding at fixed physical parameters.
Replacing the trigger to erase this discrepancy would be a new closure.

Q.A proposes qualifying a quasi-steady mean-service embedding while preserving
the inactive criterion. Q.B would retain this native law and explicitly amend
the compact admission criterion to a predictive-error diagnostic. Neither
option is selected. A source-family prediction difference is not a pure
state-removal theorem.

\newpage
\section*{3. What can be proved about recovery and the TST tail}
For the source acceleration on level terrain, suppose $a_0>0$ and
$v_{\max}>u>0$, where $u$ is free speed. While realized speed $v\le u$,
\begin{equation}
a(v)=a_0(1-v/v_{\max})\ge a_{\min}
=a_0(1-u/v_{\max})>0,\qquad
T_{\rm rec}\le (u-v_0)_+/a_{\min}.
\end{equation}
The recovery bound applies \emph{after downstream flow permits unrestricted
acceleration}, with initial desired speed at least $v_0$. A discrete source
update may add one time bin. Downstream congestion can postpone that start;
finite descriptor count alone does not prove eventual release. These stated
conditions are sufficient analytic conditions, not new adopted geometry.

After the last offered arrival, total inventory $M(t)$ is nonnegative and
nonincreasing by conservation. Suppose the eventual joined law independently
certifies constants $H>0$ and $0<\eta\le1$ for all $t\ge T$:
\begin{equation}
M(t+H)\le(1-\eta)M(t).
\end{equation}
Then $M(T+jH)\le(1-\eta)^j M(T)$. Bounding each interval by its start inventory
and summing the geometric series yields
\begin{equation}
\int_T^\infty M(t)\,dt
\le H\sum_{j=0}^{\infty}M(T+jH)
\le \frac{H M(T)}{\eta}.
\end{equation}
The right side has vehicle-equivalent seconds as units. It certifies a residual
tail for complete TST without declaring a nonempty road cleared. The constants
must be established from the joined law, not assumed from positive permission,
finite population, fitted terminal values or the absence of newly arriving traffic.

\textbf{Remaining obligations.} The adapted reference's simultaneous-event
coupling and whole draining certificate are not proved. The compact mean map
still needs its own inputs, regimes and consistent coupled solution. Common
rate-to-exit/TST leverage, policy dominance, state necessity and useful gradients
remain unanswered. The current result is new interface/compact qualification
evidence and a concrete S3-G2-Q discussion, not Stage 3 closure.

\section*{Source locations and status}
Supplied Laval--Daganzo (2006), Sections 2.2 and Appendix A; Daganzo--Laval
RR-2003-2, Sections 2--4; RM0010 (2013), Sections 4 and 6,
\url{https://arxiv.org/pdf/1310.2660}. Node and moment wrappers here are declared
review proposals. The tail inequality and inactive hand example are reviewer
analysis. See the adjacent evidence and assessment Markdown records for
source-access limits and the two candidate contracts.
\end{document}
